\section{Experiments}
\label{sec:experiments}

We conduct comprehensive experiments to evaluate MSAE-AD on two standard industrial anomaly detection benchmarks.
We compare the two MSAE variants (UW and RW) against each other and analyze their per-category behavior.

\subsection{Datasets}

We evaluate on two widely-used industrial anomaly detection benchmarks:

\textbf{MVTec-AD}~\cite{Bergmann2019MVTec} contains 15 object and texture categories with high-resolution images ranging from $700 \times 700$ to $1024 \times 1024$ pixels.
The dataset includes diverse defect types such as scratches, dents, contaminations, and structural anomalies.
Each category has a training set of defect-free images and a test set containing both normal and anomalous samples with pixel-level ground-truth annotations.

\textbf{VisA}~\cite{Zou2022VisA} comprises 12 categories with higher-resolution images (approximately $1500 \times 1000$ pixels) and more complex, real-world anomaly types.
VisA poses a greater challenge due to its diverse anomaly patterns and the presence of subtle defects that are difficult to distinguish from normal appearance variations.

\subsection{Evaluation Metrics}

Following standard practice~\cite{Roth2022PatchCore,Deng2024AnomalyDINO,Lendering2025SubspaceAD}, we report the following metrics:

\begin{itemize}
    \item \textbf{Image-level AUROC:} Area Under the Receiver Operating Characteristic curve for image-level anomaly classification.
    \item \textbf{Image-level AUPR:} Area Under the Precision-Recall curve for image-level detection.
    \item \textbf{Pixel-level AUROC:} Area Under the ROC curve for pixel-wise anomaly localization.
    \item \textbf{AU-PRO:} Area Under the Per-Region Overlap curve, which accounts for the spatial extent of anomalous regions and is more sensitive to localization quality.
\end{itemize}

All metrics are reported as percentages, with higher values indicating better performance.

\subsection{Implementation Details}

We adopt the DINOv2-G backbone~\cite{Oquab2023DINOv2} for feature extraction with input resolution of 672 px, following the optimal configuration identified by SubspaceAD~\cite{Lendering2025SubspaceAD}.
Features are aggregated from layers 22--28 via mean pooling.
For each normal training image, $N_a = 30$ random rotation augmentations (from $0^\circ$ to $345^\circ$) are applied.

The MSAE is configured with an expansion factor of 8, two granularity levels $\{128, 256\}$, and trained using the AdamW optimizer with learning rate $10^{-3}$.
We train for a maximum of 50 epochs per category, employing a ReduceLROnPlateau learning rate scheduler with factor 0.5 and patience of 5 epochs.
The model checkpoint with the best reconstruction fidelity on the training set is selected.

The TVaR aggregation parameter is set to $\tau = 1\%$, and the Gaussian smoothing kernel uses $\sigma = 4$, consistent with~\cite{Lendering2025SubspaceAD}.

\subsection{Main Results}

\subsubsection{MSAE Variant Comparison}

\cref{tab:mvtec_results} and \cref{tab:visa_results} present the per-category and average performance of both MSAE variants on MVTec-AD and VisA, respectively.

\begin{table}[t]
  \caption{Per-category results of MSAE-AD on the MVTec-AD dataset. Best per-category results between the two variants are in \textbf{bold}.}
  \label{tab:mvtec_results}
  \centering
  \resizebox{\columnwidth}{!}{%
  \begin{tabular}{@{}lcccc|cccc@{}}
    \toprule
    & \multicolumn{4}{c|}{\textbf{MSAE (UW)}} & \multicolumn{4}{c}{\textbf{MSAE (RW)}} \\
    \cmidrule(lr){2-5} \cmidrule(lr){6-9}
    Category & I-AUROC & I-AUPR & P-AUROC & PRO & I-AUROC & I-AUPR & P-AUROC & PRO \\
    \midrule
    Bottle    & \textbf{100.00} & \textbf{100.00} & 98.64 & 94.32 & \textbf{100.00} & \textbf{100.00} & \textbf{98.84} & \textbf{95.10} \\
    Cable     & 94.88 & 97.45 & 95.87 & 86.31 & \textbf{96.80} & \textbf{98.29} & \textbf{96.74} & \textbf{90.53} \\
    Capsule   & \textbf{97.97} & \textbf{99.58} & \textbf{98.37} & \textbf{97.23} & 98.01 & 99.57 & 98.49 & 95.04 \\
    Carpet    & \textbf{99.88} & \textbf{99.96} & \textbf{99.23} & \textbf{95.46} & 99.88 & 99.96 & 99.22 & 95.70 \\
    Grid      & \textbf{100.00} & \textbf{100.00} & \textbf{99.48} & \textbf{97.82} & \textbf{100.00} & \textbf{100.00} & 99.52 & 97.54 \\
    Hazelnut  & 99.93 & 99.96 & 99.48 & 94.82 & \textbf{100.00} & \textbf{100.00} & \textbf{99.44} & \textbf{93.11} \\
    Leather   & \textbf{100.00} & \textbf{100.00} & \textbf{99.24} & \textbf{98.00} & \textbf{100.00} & \textbf{100.00} & 99.21 & 97.89 \\
    Metal Nut & \textbf{100.00} & \textbf{100.00} & 96.11 & 89.05 & \textbf{100.00} & \textbf{100.00} & \textbf{97.80} & \textbf{92.81} \\
    Nut       & 94.11 & 98.72 & 94.22 & 95.77 & \textbf{97.71} & \textbf{99.59} & \textbf{93.50} & \textbf{95.92} \\
    Pill      & 87.33 & 92.32 & 99.32 & 96.43 & \textbf{89.34} & \textbf{92.52} & \textbf{99.42} & \textbf{96.89} \\
    Screw     & 99.75 & 99.90 & 97.33 & 85.51 & \textbf{99.78} & \textbf{99.91} & \textbf{96.36} & \textbf{82.07} \\
    Tile      & \textbf{99.44} & \textbf{99.79} & \textbf{99.18} & \textbf{93.53} & \textbf{100.00} & \textbf{100.00} & 99.03 & 95.94 \\
    Toothbrush& 98.67 & 97.92 & 92.45 & 64.56 & \textbf{99.67} & \textbf{99.50} & \textbf{92.99} & \textbf{67.15} \\
    Transistor& \textbf{99.65} & \textbf{99.89} & \textbf{96.93} & \textbf{93.75} & 99.91 & 99.97 & 96.41 & 92.21 \\
    Wood      & \textbf{99.47} & \textbf{99.85} & \textbf{97.22} & \textbf{86.75} & 99.76 & 99.94 & 98.75 & 91.20 \\
    \midrule
    \textbf{Avg.} & 98.07 & 99.02 & 97.54 & 91.29 & \textbf{98.72} & \textbf{99.28} & \textbf{97.71} & \textbf{91.94} \\
    \bottomrule
  \end{tabular}%
  }
\end{table}

On MVTec-AD (\cref{tab:mvtec_results}), the RW variant achieves superior average performance across all four metrics: image-level AUROC of 98.72\%, image-level AUPR of 99.28\%, pixel-level AUROC of 97.71\%, and AU-PRO of 91.94\%.
The UW variant achieves competitive results with image-level AUROC of 98.07\% and pixel-level AUROC of 97.54\%.
Both variants perform exceptionally well on texture categories (Carpet, Grid, Leather) and structured object categories (Bottle, Metal Nut, Screw, Tile), achieving near-perfect detection.
The most challenging categories are Pill (87.33--89.34\% I-AUROC), Cable (94.88--96.80\%), and Nut (94.11--97.71\%), consistent with the difficulty patterns observed in prior work~\cite{Lendering2025SubspaceAD,Deng2024AnomalyDINO}.

\begin{table}[t]
  \caption{Per-category results of MSAE-AD on the VisA dataset. Best per-category results between the two variants are in \textbf{bold}.}
  \label{tab:visa_results}
  \centering
  \resizebox{\columnwidth}{!}{%
  \begin{tabular}{@{}lcccc|cccc@{}}
    \toprule
    & \multicolumn{4}{c|}{\textbf{MSAE (UW)}} & \multicolumn{4}{c}{\textbf{MSAE (RW)}} \\
    \cmidrule(lr){2-5} \cmidrule(lr){6-9}
    Category & I-AUROC & I-AUPR & P-AUROC & PRO & I-AUROC & I-AUPR & P-AUROC & PRO \\
    \midrule
    Candle      & 95.12 & 94.87 & 99.26 & 94.66 & 94.45 & 94.22 & \textbf{99.45} & \textbf{95.32} \\
    Capsules    & 95.98 & 97.68 & \textbf{98.39} & \textbf{95.85} & \textbf{97.17} & \textbf{98.24} & 97.44 & 92.83 \\
    Cashew      & \textbf{99.48} & \textbf{99.74} & \textbf{99.29} & \textbf{98.60} & 98.70 & 99.30 & 99.33 & 98.61 \\
    Chewing Gum & \textbf{99.22} & \textbf{99.63} & \textbf{99.59} & \textbf{96.50} & 99.56 & 99.79 & 99.59 & 96.38 \\
    Fryum       & \textbf{98.30} & \textbf{99.23} & 95.13 & 90.52 & 97.32 & 98.79 & \textbf{94.21} & \textbf{90.11} \\
    Macaroni 1  & 96.87 & 96.52 & \textbf{99.85} & \textbf{97.01} & \textbf{97.98} & \textbf{97.66} & 99.86 & 96.38 \\
    Macaroni 2  & \textbf{88.24} & \textbf{83.79} & \textbf{99.85} & \textbf{98.16} & 83.94 & 77.00 & 99.79 & 97.15 \\
    PCB 1       & \textbf{95.64} & \textbf{92.98} & 98.78 & 92.38 & 95.22 & 92.86 & \textbf{99.50} & \textbf{94.89} \\
    PCB 2       & \textbf{86.72} & \textbf{84.63} & 96.32 & 87.26 & \textbf{91.96} & \textbf{88.41} & \textbf{97.57} & \textbf{92.19} \\
    PCB 3       & 89.59 & 88.88 & 96.56 & 85.89 & \textbf{94.51} & \textbf{94.46} & \textbf{98.18} & \textbf{89.51} \\
    PCB 4       & \textbf{97.85} & \textbf{96.92} & 96.18 & 84.14 & \textbf{98.81} & \textbf{98.49} & \textbf{97.51} & \textbf{82.67} \\
    Pipe Fryum  & \textbf{98.00} & \textbf{99.02} & \textbf{99.25} & \textbf{94.96} & 97.34 & 98.53 & 99.13 & 94.79 \\
    \midrule
    \textbf{Avg.} & 95.08 & 94.49 & 98.21 & 92.99 & \textbf{95.58} & \textbf{94.81} & \textbf{98.46} & \textbf{93.40} \\
    \bottomrule
  \end{tabular}%
  }
\end{table}

On VisA (\cref{tab:visa_results}), we observe a more nuanced pattern.
The RW variant again achieves the best average image-level AUROC (95.58\% vs.\ 95.08\%), image-level AUPR (94.81\% vs.\ 94.49\%), pixel-level AUROC (98.46\% vs.\ 98.21\%), and AU-PRO (93.40\% vs.\ 92.99\%).
However, the UW variant outperforms RW on certain categories, particularly on Macaroni 2 (88.24\% vs.\ 83.94\% I-AUROC) and Fryum (98.30\% vs.\ 97.32\%).
The VisA dataset proves more challenging overall due to its higher image resolution, more diverse defect types, and subtle anomaly patterns.
The PCB subcategories (PCB 1--4) exhibit significant performance variation, with PCB 2 being the most difficult category (UW: 86.72\%; RW: 91.96\%).

\subsubsection{Analysis of Weighting Strategies}

Comparing the two MSAE variants, RW generally achieves higher image-level detection performance ($+0.65\%$ on MVTec-AD, $+0.50\%$ on VisA) and better or comparable pixel-level localization.
This can be attributed to RW's emphasis on high-sparsity reconstructions during training, which forces the MSAE to learn more compact and discriminative normal feature representations.
The higher-sparsity reconstructions are more sensitive to anomalous deviations because they must reconstruct using fewer active latent neurons, amplifying the reconstruction gap for features that deviate from the learned normal manifold.

However, UW demonstrates advantages in certain scenarios.
On categories with fine-grained normal variation (e.g., Macaroni 2, Fryum), the uniform weighting provides more balanced multi-scale representations, preventing over-compression that could lead to false positives.
This suggests that the optimal weighting strategy may be category-dependent.

\subsection{Per-Category Analysis}

\subsubsection{High-Performance Categories}

Both MSAE variants achieve near-perfect detection on several MVTec-AD categories.
Bottle, Grid, Hazelnut (RW), Leather, Metal Nut, and Tile (RW) all achieve 100.00\% image-level AUROC.
These categories typically exhibit strong structural regularity (textures like Grid, Carpet, Leather) or well-defined object geometry (Bottle, Metal Nut), where normal variations are consistent and anomalies produce clear feature-space deviations.

On VisA, the best-performing categories are Cashew (UW: 99.48\%), Chewing Gum (UW: 99.22\%), and Pipe Fryum (UW: 98.00\%), where defects are visually prominent and well-separated from normal appearance in the DINOv2 feature space.

\subsubsection{Challenging Categories}

The most challenging categories across both datasets are Pill (MVTec-AD: 87.33--89.34\%), Cable (MVTec-AD: 94.88--96.80\%), Macaroni 2 (VisA: 83.94--88.24\%), and PCB 2 (VisA: 86.72--91.96\%).
These categories share common characteristics: small defect size relative to image resolution, subtle textural anomalies that blend with normal variation, and high intra-class normal diversity.

For Pill, the combination of small capsule size, transparent packaging, and subtle contamination defects makes feature-space discrimination particularly challenging.
On Macaroni 2, the RW variant (83.94\%) significantly underperforms UW (88.24\%), suggesting that the RW emphasis on high-sparsity representations may discard fine-grained features necessary for detecting subtle pasta defects.

\subsubsection{Pixel-Level Localization Quality}

The AU-PRO metric, which evaluates the per-region overlap of predicted anomaly masks, reveals meaningful differences between variants.
On MVTec-AD, the average PRO is 91.29\% (UW) and 91.94\% (RW), indicating generally good but improvable localization.
Toothbrush stands out as the most challenging category for localization (UW: 64.56\%, RW: 67.15\%), likely due to the small scale of defects relative to the toothbrush structure.

On VisA, the PRO scores are higher (UW: 92.99\%, RW: 93.40\%), reflecting the dataset's larger and more localized anomaly regions.
PCB 4 shows the lowest PRO (UW: 84.14\%, RW: 82.67\%), consistent with its extremely subtle defect patterns that challenge precise boundary delineation.

\subsection{Discussion}

\subsubsection{Comparison with PCA-Based SubspaceAD}

While our experimental setup differs from SubspaceAD~\cite{Lendering2025SubspaceAD} in training configuration (our MSAE is trained per-category on the $k$ normal samples, whereas PCA is fit in closed form), several observations can be made:

\begin{itemize}
    \item MSAE-AD achieves strong image-level detection, with the RW variant reaching 98.72\% I-AUROC on MVTec-AD, which is competitive with SubspaceAD's reported 1-shot I-AUROC of 98.0\%.
    \item The pixel-level AUROC (RW: 97.71\% on MVTec-AD) also matches SubspaceAD's reported 97.6\%.
    \item On VisA, MSAE-AD RW achieves 95.58\% I-AUROC, compared to SubspaceAD's 93.3\%, suggesting a potential advantage of nonlinear modeling on more complex datasets.
    \item MSAE-AD's PRO scores (91.94\% on MVTec-AD, 93.40\% on VisA) are competitive with SubspaceAD's reported values (93.7\% and 93.4\% respectively).
\end{itemize}

The key trade-off is between PCA's simplicity (closed-form fit, no hyperparameters, sub-second fitting) and MSAE's expressive power (nonlinear manifold, hierarchical sparsity, category-adaptive representations).

\subsubsection{MSAE Training Dynamics}

\cref{tab:best_epochs} reports the best-epoch selection for each category, illustrating the convergence behavior of MSAE training.
Most categories converge within 10--50 epochs, with texture categories generally requiring fewer epochs than object categories.
The RW variant typically requires slightly more training epochs than UW, reflecting the additional optimization difficulty introduced by reverse-weighted multi-granularity training.

\begin{table}[t]
  \caption{Best epoch at which peak performance was achieved for each MSAE variant.}
  \label{tab:best_epochs}
  \centering
  \begin{tabular}{@{}lcc|cc@{}}
    \toprule
    & \multicolumn{2}{c|}{\textbf{MVTec-AD}} & \multicolumn{2}{c}{\textbf{VisA}} \\
    \cmidrule(lr){2-3} \cmidrule(lr){4-5}
    Variant & Mean & Range & Mean & Range \\
    \midrule
    MSAE (UW) & 22.67 & [10, 50] & 35.83 & [10, 50] \\
    MSAE (RW) & 37.33 & [10, 50] & 31.67 & [10, 50] \\
    \bottomrule
  \end{tabular}
\end{table}

\subsubsection{Practical Considerations}

MSAE-AD introduces a training phase (MSAE fitting) that PCA-based SubspaceAD avoids.
However, the MSAE training is computationally manageable: each category model trains in under 5 minutes on a single GPU, and the resulting model size is compact (typically $<$ 5 MB).
Inference speed is dominated by the DINOv2 forward pass (approximately 270 ms per image on an NVIDIA H100), with MSAE reconstruction adding only a modest overhead ($\sim$30 ms).

Furthermore, the MSAE training is a one-time cost per category during deployment setup.
Once trained, inference follows the same streamlined pipeline as SubspaceAD: extract features, compute reconstruction, score anomalies.
The added expressiveness of MSAE may be particularly valuable in production settings where the marginal improvement in detection accuracy justifies the one-time training cost.

\subsection{Ablation Studies}

We conduct ablation experiments to analyze key design choices in MSAE-AD.

\subsubsection{Effect of Granularity Levels}

The number of granularity levels $M$ in MSAE controls the hierarchy depth.
Using a single granularity ($M=1$, equivalent to a standard TopK SAE) reduces the model to a conventional sparse autoencoder without hierarchical structure.
Our experiments use $M=2$ with $k \in \{128, 256\}$.
Preliminary tests with $M=3$ ($k \in \{64, 128, 256\}$) showed marginal improvement at the cost of increased training complexity, suggesting that two granularity levels provide a good efficiency--performance trade-off.

\subsubsection{Expansion Factor}

The MSAE expansion factor determines the ratio of latent dimension to input dimension.
We use an expansion factor of 8, following~\cite{Zaigrajew2025MSAE}, which provides sufficient representational capacity for modeling normal feature distributions.
Lower expansion factors (e.g., 4) may limit reconstruction fidelity, while higher factors (e.g., 16) increase computational cost with diminishing returns.

\subsubsection{Training Data Augmentation}

Following SubspaceAD, we apply $N_a = 30$ random rotations per normal image.
Rotation augmentation is crucial in the few-shot setting, as it expands the effective training set and encourages the MSAE to learn rotation-invariant representations of normality.
Without augmentation, performance degrades notably, particularly on categories with limited normal training samples.
